\chapter{Sledljivost programske opreme in rezultatov}
\label{app:sledljivost}

Končna evalvacija ima identifikator
\texttt{STM32N6-HAZARD6-FINAL-003}, poskus pa oznako \texttt{FE4-A01}.
Manifest ima naslednjo kontrolno vsoto
SHA-256:
\begin{center}
  \small\texttt{0355fca0f9a3d26be1e41b4fd9896f7}\newline
  \texttt{fb7c9f567557a7dde23826dbb36c1342b}.
\end{center}

Glavni rezultati in grafi so shranjeni v naslednjih imenikih:
\begin{itemize}
  \item \path{experiments/results/hazard6_final_evaluation_v4_independent},
  \item \path{experiments/results/deployment_footprint}.
\end{itemize}
Fotografije fizične izvedbe skupaj s kontrolnimi vsotami izvirnikov so
shranjene v imeniku \path{Doc/thesis_photos}.

Postopek priprave kratkotrajnih dogodkov je v programu
\path{ml/prepare_transient_event_windows.py}, povečanje učnega nabora pa v
\path{ml/prepare_patch_balanced_augmentation.py}. Končna učna nastavitev je v
datoteki \path{ml/configs/hazard5v5_other_yamnet1024_tqe.yaml}, zagon učenja pa
je dokumentiran v programu \path{ml/train_hazard5v5_other_yamnet1024.ps1}.
Datoteke z rezultati hranijo tudi razrešeno učno nastavitev, metrike po prehodih
in kontrolne vsote modelov.

Izvorna koda, analizni programi in zapisi o razvoju so v javnem repozitoriju
\url{https://github.com/aleksMuhicFri2/stm32n6-edge-audio-classifier}. Zaradi
licenčnih in velikostnih omejitev javne zbirke zvoka niso vključene v
repozitorij; manifesti hranijo izvorne oznake in kontrolne vsote uporabljenih
dražljajev.
